\section*{Overview}

These notes are the part of the competitive programming archive that is not tied to one contest year.
The goal is simple: keep the ideas that I actually want to reuse in a form that is still readable after
the adrenaline of the original contest is gone.

I care about three things here:

\begin{itemize}
  \item the structural reason an algorithm works,
  \item the shape of the implementation I would trust in contest code,
  \item the failure modes that usually cost time or wrong answers.
\end{itemize}

\section*{How the Notes Are Written}

Each topic keeps \texttt{note.tex} as the written source of truth and \texttt{code.cpp} as the implementation
source of truth. The site renders the TeX into HTML and shows the C++ file separately, so the prose can stay
focused while the code remains directly copyable.

The finished notes are meant to be public-facing study material, not private scraps. That means the emphasis is
on explanations, invariants, and tradeoffs instead of dumping templates with no context.

\section*{What ``Finished'' Means Here}

A finished note should answer the questions I usually care about when revising:

\begin{itemize}
  \item What kind of contest situations make this technique the right tool?
  \item What invariant is doing the real work?
  \item Which implementation shape is the least fragile under time pressure?
  \item Which edge cases usually break the first version?
\end{itemize}

Some topics are still only planned. They are listed openly as planned instead of pretending that a shallow page
is enough.

\section*{References Behind the Section}

The explanations are rewritten in my own words, but the topic map and many implementation choices were informed by:

\begin{itemize}
  \item \href{https://codeforces.com/catalog}{Codeforces Catalog},
  \item \href{https://cp-algorithms.com/}{cp-algorithms},
  \item \href{https://usaco.guide/}{USACO Guide},
  \item \href{https://usaco.guide/CPH.pdf}{Competitive Programmer's Handbook},
  \item \href{https://github.com/kth-competitive-programming/kactl}{KACTL},
  \item \href{https://oi-wiki.org/}{OI Wiki}.
\end{itemize}
